\documentclass[11pt,a4paper]{article}
\usepackage[T1]{fontenc}
\usepackage[ngerman]{babel}
\usepackage{lmodern,microtype}
\usepackage[margin=23mm,headheight=15pt]{geometry}
\usepackage{amsmath,amssymb,mathtools}
\usepackage{booktabs,longtable,tabularx,array}
\usepackage{xcolor,tikz}
\usetikzlibrary{arrows.meta,positioning,calc,fit,backgrounds}
\usepackage{enumitem,xurl,hyperref,fancyhdr}
\definecolor{navy}{HTML}{16324F}
\definecolor{teal}{HTML}{087F8C}
\definecolor{amber}{HTML}{A85D08}
\definecolor{pale}{HTML}{F0F5F8}
\definecolor{red}{HTML}{9C3242}
\hypersetup{colorlinks=true,linkcolor=navy,urlcolor=teal,citecolor=teal,pdftitle={TFPT: Die gemeinsame Quelle und der Weg zur Gesamtlösung},pdfauthor={Sitzungssynthese für Stefan Hamann}}
\pagestyle{fancy}\fancyhf{}\fancyhead[L]{\small TFPT \textbar\ Gesamtsynthese der Sitzung}\fancyhead[R]{\small Stand: 30. September 2026}\fancyfoot[C]{\thepage}
\setlength{\parindent}{0pt}\setlength{\parskip}{5pt plus 1pt}
\setlength{\emergencystretch}{3em}
\setlist{itemsep=3pt,topsep=4pt,leftmargin=*}
\newcolumntype{Y}{>{\raggedright\arraybackslash}X}
\newcolumntype{L}[1]{>{\raggedright\arraybackslash}p{#1}}
\newcommand{\E}{\textcolor{teal}{\textbf{E}}}
\newcommand{\N}{\textcolor{teal}{\textbf{N}}}
\newcommand{\B}{\textcolor{amber}{\textbf{B}}}
\newcommand{\Opend}{\textcolor{amber}{\textbf{O}}}
\newcommand{\Xclaim}{\textcolor{red}{\textbf{X}}}
\newcommand{\C}{\mathbb C}\newcommand{\R}{\mathbb R}\newcommand{\Z}{\mathbb Z}
\newcommand{\tr}{\operatorname{tr}}\newcommand{\diag}{\operatorname{diag}}
\newcommand{\Tr}{\operatorname{Tr}}
\newcommand{\rank}{\operatorname{rank}}\newcommand{\Span}{\operatorname{span}}
\newcommand{\Id}{\mathrm{id}}\newcommand{\ket}[1]{|#1\rangle}\newcommand{\bra}[1]{\langle#1|}
\newcommand{\abs}[1]{\left|#1\right|}\newcommand{\norm}[1]{\left\|#1\right\|}
\newcommand{\file}[1]{\texttt{\nolinkurl{#1}}}
\newcommand{\key}[1]{\par\smallskip\noindent\colorbox{pale}{\parbox{\dimexpr\linewidth-2\fboxsep\relax}{{\color{navy}#1}}}\par\smallskip}
\newcommand{\picturetext}[1]{\textit{Bildlich:} #1}
\tikzset{box/.style={draw=navy,rounded corners=3pt,fill=pale,align=center,font=\small,inner sep=7pt},arr/.style={-{Latex[length=2mm]},thick,draw=teal},dasharr/.style={arr,dashed}}

